\chapter{Conclusion}
\label{ch:conclusion}

\section{Summary}

This thesis asked whether semantic entropy provides a better
hallucination-risk and abstention signal than surface-form uncertainty,
and how that answer depends on the surface-form variability of the
task, the clustering backend, and the scale of the generator. The
method follows Farquhar et al.~\cite{semantic_entropy_nature_2024};
the contribution here is a reproducible reimplementation evaluated on
a full cross-seed grid --- three open generators, three short-answer
datasets and a long-form biography task, three decoding seeds under
the original paper's sampling and labelling protocol --- plus a
claim-level long-form extension, rather than a new semantic-entropy
method.

The answer is regime-dependent, in both directions and more sharply
than expected. On short-answer tasks every sampling-based estimator
detects errors well (AUROC 0.74--0.88), but entropy over normalized
strings --- the cheapest score in the family --- is the strongest
single method in six of nine model--dataset cells, and no semantic,
kernel, probe, or self-evaluation method beats it significantly
anywhere. A record-level analysis explains why: dataset normalization
already removes almost all surface variation before any semantic step
runs, so clustering has little left to merge. On long-form biography the
same surface scores collapse to chance while every detector that escapes
raw string-counting recovers a statistically significant signal for
every generator; the strongest long-form scores come from
probability-weighted naive entropy on the two smaller models
(AUROC 0.77--0.78) and from the single-pass P(True) on the largest
(0.78). Under the paper-faithful protocol, the short-answer
semantic-entropy averages land in the same band as the original work's
reported 0.790, with the same method ordering.

Two further findings qualify the method's practical use. First, the
clustering backend is result-determining on long-form inputs: across
NLI cross-encoders and an LLM judge, backends span chance-level to
clearly useful behavior on the same saved samples, and the same
artifact reappears at claim level, where exact-match clustering wrongly
makes the largest model look worse than chance. Second, the claim-level FactualBio extension shows that
semantic entropy's discriminative power decreases monotonically with
generator scale (AUROC 0.688 $\rightarrow$ 0.644 $\rightarrow$ 0.521
from 7B to 70B): larger models repeat their confabulations more
consistently, eroding exactly the inconsistency that the method
measures. Hidden-state probes invert this picture. On short answers
they are merely a cheap alternative --- strictly weaker than
sampling-based entropy, and converging to the same performance whether
trained on correctness or on the semantic-entropy threshold --- but on
long-form biography the SEP probe is competitive and, unlike every
sampling-based score, its accuracy \emph{rises} with model scale ---
though naive entropy, and on the 70B model P(True), score higher. Where
a larger model's growing self-consistency erodes entropy- and
agreement-based detection, a probe that reads the model's own
representation does not, making it the natural complement to semantic
entropy precisely where entropy is weakest.

\section{Contributions}

The substantive method comes from Farquhar et
al.~\cite{semantic_entropy_nature_2024}. The thesis contributes six
main artifacts:

\begin{enumerate}
  \item a reproducible \texttt{semantic\_entropy} implementation for
        normalization, clustering under multiple entailment backends,
        the four-estimator entropy family, correctness labeling, and
        selective-prediction metrics;
  \item a complete cross-seed experiment record --- three generators,
        four datasets, three seeds under paper-faithful sampling and
        accuracy labelling --- with Slurm templates, JSONL outputs,
        per-run and cross-seed metric summaries, and record-level
        paired-bootstrap significance kept separate from seed
        stability;
  \item extension baselines covering P(True), in- and
        out-of-distribution embedding regression, Kernel Language
        Entropy under two encoders, and the Semantic Entropy
        Probe~\cite{kossen_semantic_entropy_probes_2024}, with the
        probe-equivalence observation that correctness-trained and
        SE-trained probes are interchangeable in practice;
  \item a long-form evaluation at two granularities: the
        whole-paragraph biography task under the same
        baseline-versus-semantic design, and the claim-level
        FactualBio extension --- biographies generated by the models
        under test, decomposed into claims, re-queried, and scored by
        context-conditioned NLI-clustered semantic entropy with
        judge-graded labels --- which yields the scale trend and the
        clustering-artifact result;
  \item record-level worked audits drawn from the canonical runs that
        trace where clustering actually happens (normalization versus
        NLI), show a confident systematic error the method cannot catch,
        and explain why the surface baseline collapses on long text,
        providing mechanistic rather than purely metric-level evidence
        for the main conclusions (Section~\ref{sec:hand-audits});
  \item a sampling-temperature sweep ($T\in\{0.1,0.3,0.5,0.7\}$ against
        the canonical $T=1.0$) over the full short-answer and biography
        grid, with the numbers recomputed from raw scores as an
        independent check, isolating the opposite temperature dependence
        of the two detector families --- sampling-based entropy
        improving as the decoder warms, the hidden-state probe declining
        --- and the regime-dependent crossover this implies for
        deployment (Section~\ref{sec:temperature-sensitivity}).
\end{enumerate}

The full source code, result tables, Slurm templates, and LaTeX thesis
source are available at \url{https://github.com/D3prave/thesis}.
The Python implementation lives in the \texttt{src/semantic\_entropy/}
package; entry points are documented in \texttt{README.md}.

\section{Limitations}

The thesis is deliberately narrow in scope. The detection target is the
confabulation setting defined in Chapter~\ref{ch:background}.
High-confidence systematic errors are by construction not detected by
this method --- a limitation the FactualBio scale trend makes concrete,
since consistent confabulation at 70B is precisely the confident-error
case; Simhi et al.~\cite{simhi_choke_2025} discuss this setting.

Within the empirical results, the main residual limitations are the
automated LLM judging behind every long-form label (whole-paragraph
and claim-level alike), the three-seed basis of the stability
estimates, and the English-only, biography-only long-form scope.
Record-level issues --- alias-type oracle failures on TriviaQA,
semantically consistent refusals on biography, and SVAMP's
integer-arithmetic structure --- are documented in the audits; the
fuller discussion is in Section~\ref{sec:discussion-limitations}.

\section{Future Work}

Three extensions follow directly from the limitations above. First,
the automated claim and paragraph labels should be validated against a
human-graded
subsample, since every long-form conclusion inherits the judge's error
modes. Second, the FactualBio scale trend invites extension beyond 70B
and beyond biographical claims: if consistency-based detection
continues to erode with scale, the practical future of confabulation
detection lies with methods that do not require the model to re-roll
its errors, such as hidden-state probes or retrieval-grounded
verification. Third, the long-form evaluation should extend beyond
English single-reference biography and be paired with a dedicated
detector for the semantically consistent refusals that low entropy
cannot flag.
